\documentclass[11pt]{article}
\usepackage[margin=1in]{geometry}
\usepackage{amsmath,amssymb,amsthm,hyperref}
\hypersetup{colorlinks=true,linkcolor=blue,urlcolor=blue}
\newcommand{\Q}{\mathbb Q}
\newcommand{\sech}{\operatorname{sech}}
\newcommand{\csch}{\operatorname{csch}}
\newcommand{\Res}{\operatorname{Res}}
\newtheorem{theorem}{Theorem}
\title{The coefficient relations for squared Berndt integrals\\MathDB 357280}
\author{Research note}
\date{17 September 2026}
\begin{document}
\maketitle

\section*{Task 1. Prove both coefficient relations}
Put $G=\Gamma(1/4)$. For every integer $p\ge1$, the two evaluations
in Xu--Zhao \cite[Theorems 6.3 and 8.12]{XZ} have the form
\begin{align}
 I_p^-:=\int_0^\infty\frac{x^{4p+1}\,dx}{(\cos x-\cosh x)^2}
   &=i_p\frac{G^{8p}}{\pi^{2p}}
                  +j_p\frac{G^{8p+8}}{\pi^{2p+4}},\label{minus}\\
 I_p^+:=\int_0^\infty\frac{x^{4p+1}\,dx}{(\cos x+\cosh x)^2}
   &=g_p\frac{G^{8p}}{\pi^{2p}}
                  +h_p\frac{G^{8p+8}}{\pi^{2p+4}},\label{plus}
\end{align}
with rational coefficients. This specifies the theorem references
omitted from the MathDB transcription. We prove Conjecture 8.14
of the current source, without changing its range or conclusion.

\begin{theorem}
For every positive integer $p$,
\[
 \boxed{\displaystyle
 g_p=-\frac{2^{2p-1}-(-1)^p}{2^{2p-1}}\,i_p,
 \qquad
 h_p=-\frac{2^{2p-1}+(-1)^p}{2^{2p-1}}\,j_p.}
\]
\end{theorem}

The proof also constructs all four rational coefficients. The
essential connection between the two integral families is the
degree-two Eisenstein identity, followed by a modular transformation
from $i$ to $(1+i)/2$.

\subsection*{1. A common residue calculation}
Fix $p\ge1$ and write $w=4p$, $a=(-1)^p2^{2p-1}$. For $y>0$ set
\[
 S_-(y)=\sum_{n\ge1}n^w\csch^2(ny),\qquad
 S_+(y)=\sum_{\substack{k\ge1\\k\text{ odd}}}
                         k^w\sech^2(ky/2).
\]
We first establish
\begin{align}
 I_p^-&=a\pi^{w+1}\{(w+1)S_-(\pi)+\pi S_-'(\pi)\},\label{res-minus}\\
 I_p^+&=-\frac{a}{2^w}\pi^{w+1}
                    \{(w+1)S_+(\pi)+\pi S_+'(\pi)\}.
                    \label{res-plus}
\end{align}
These formulas can also be read from the source's contour results;
we give the details to fix their relative sign and scale.

Let $a_0=(1+i)/2$, $b_0=(1-i)/2$. The factorizations are
\[
 \cos z+\cosh z=2\cosh(a_0z)\cosh(b_0z),\qquad
 \cos z-\cosh z=-2\sinh(a_0z)\sinh(b_0z).
\]
Use positively oriented triangles with vertices $0,R,iR$, taking
$R=2\pi N$ for the plus sign and $R=(2N+1)\pi$ for the minus sign.
The first hyperbolic factor on the oblique edge has modulus at least
one, and the second has modulus at least $\sinh(R/2)$. The oblique
integral is $O_p(R^{w+2}e^{-R})$ and tends to zero. The two axial
integrals have the same sign and sum to twice the real integral.
For the minus sign the apparent singularity at zero is removable:
the quotient is $O(z^{4p-3})$ there. Thus for either sign
the integral equals $\pi i$ times the sum of its first-quadrant residues.

For the minus sign the poles are $(1+i)t$, $t=n\pi$, $n\ge1$.
Writing $z=(1+i)t+v$, the denominator squared is
$4\sinh^2(a_0v)\sinh^2(t+b_0v)$. Since
$b_0(1+i)=1$ and $(1+i)^w=(-4)^p$, its residue is
\[
 -ia\{(w+1)t^w\csch^2t+t^{w+1}(\csch^2)'(t)\}.
\]
For the plus sign $t=(n+1/2)\pi$, $n\ge0$, the local denominator
squared is $-4\sinh^2(a_0v)\cosh^2(t+b_0v)$, and the residue is
\[
 ia\{(w+1)t^w\sech^2t+t^{w+1}(\sech^2)'(t)\}.
\]
Summing and multiplying by $\pi i$ proves
\eqref{res-minus}--\eqref{res-plus}. All sums and their differentiated
versions converge locally uniformly for $y>0$ by exponential decay.

\subsection*{2. The degree-two Eisenstein identity}
Use $q=e^{2\pi i\tau}$ and $D=(2\pi i)^{-1}d/d\tau=q,d/dq$.
Let
\[
 f(\tau)=E_w(\tau)=1-\frac{2w}{B_w}
                    \sum_{n\ge1}\sigma_{w-1}(n)q^n,
 \qquad c=-\frac{B_w}{2w}.
\]
Then $L_{w-1}(q)=\sum_{n\ge1}n^{w-1}q^n/(1-q^n)=c(f-1)$.
For $\tau=iy/\pi$, expansion of $\csch^2$ immediately gives
\begin{equation}\label{Sminus}
 S_-(y)=4c(Df)(\tau).
\end{equation}
Selecting odd $k$ in the expansion of $\sech^2$ similarly gives
\begin{equation}\label{Splus-long}
 S_+(y)=4c\{(Df)(\tau/2)
       -(2^w+4)(Df)(\tau)+2^{w+2}(Df)(2\tau)\}.
\end{equation}
For example, the selector identity before differentiation is
\[
 \sum_{k\text{ odd}}\frac{k^{w-1}q^k}{1+q^k}
 =L_{w-1}(q)-(2+2^{w-1})L_{w-1}(q^2)+2^wL_{w-1}(q^4).
\]
Its derivative gives \eqref{Splus-long}, including the chain factors.

The elementary divisor identity, with $\sigma_r(n/2)=0$ for odd $n$,
\[
 \sigma_r(2n)=(1+2^r)\sigma_r(n)-2^r\sigma_r(n/2)
\]
proves coefficient by coefficient that
\begin{equation}\label{hecke}
 f(\tau/2)+f((\tau+1)/2)+2^wf(2\tau)=(2+2^w)f(\tau).
\end{equation}
Indeed the two half-argument terms select the even Fourier indices;
the constant terms also agree. Differentiate \eqref{hecke} and
substitute in \eqref{Splus-long} to obtain the key simplification
\begin{equation}\label{Splus}
 S_+(y)=4c\{2^w(Df)(\tau)-(Df)((\tau+1)/2)\}.
\end{equation}

\subsection*{3. Two transformed derivatives}
Write $f_r=(D^rf)(i)$ for $r=0,1,2$, and put
\[
 K=(-1)^p2^{2p}=(1+i)^w,\qquad t=\frac{2K}{2^w}
                                        =\frac{(-1)^p}{2^{2p-1}}.
\]
The transformation $\gamma\tau=\tau/(\tau+1)$ is in
$\mathrm{SL}_2(\mathbb Z)$ and sends $i$ to $(1+i)/2$.
For any modular form $f$ of weight $w$, differentiating
$f(\gamma\tau)=(\tau+1)^wf(\tau)$ once and twice gives, with
$\lambda=(2\pi i)^{-1}$,
\begin{align*}
 (Df)(\gamma\tau)
   &=(\tau+1)^{w+2}Df(\tau)
                   +w\lambda(\tau+1)^{w+1}f(\tau),\\
 (D^2f)(\gamma\tau)
   &=(\tau+1)^{w+4}D^2f(\tau)
     +2(w+1)\lambda(\tau+1)^{w+3}Df(\tau)\\
   &\hspace{12mm}+w(w+1)\lambda^2(\tau+1)^{w+2}f(\tau).
\end{align*}
Differentiating $f(-1/\tau)=\tau^wf(\tau)$ at $\tau=i$ gives
$f_1=wf_0/(4\pi)$. Substituting this and $(1+i)^2=2i$ above yields
\begin{equation}\label{transformed}
 (Df)((1+i)/2)=2Kf_1,\qquad
 (D^2f)((1+i)/2)
       =K\left(-4f_2+\frac{w(w+1)}{2\pi^2}f_0\right).
\end{equation}
Now $d/dy=-2D$ for the argument $\tau=iy/\pi$, whereas it is
$-D$ for $(\tau+1)/2$. Equations \eqref{Sminus}, \eqref{Splus},
and \eqref{transformed} give
\begin{align*}
 S_-(\pi)&=4cf_1,& S_-'(\pi)&=-8cf_2,\\
 S_+(\pi)&=4c(2^w-2K)f_1,&
 S_+'(\pi)&=-8c\left((2^w+2K)f_2
                           -\frac{Kw(w+1)}{4\pi^2}f_0\right).
\end{align*}
Using $f_1=wf_0/(4\pi)$ in the residue formulas, the extra
$Kf_0$ terms cancel. We obtain
\begin{align}
 I_p^-&=4ca\pi^{w+1}\{(w+1)f_1-2\pi f_2\},\label{Iminus}\\
 I_p^+&=-4ca\pi^{w+1}\{(w+1)f_1-2\pi(1+t)f_2\}.
                                                        \label{Iplus}
\end{align}

\subsection*{4. Separate the two period powers}
Set $u=G^4/(4\pi^3)$. The classical values are
\[
 E_2(i)=3/\pi,\qquad E_4(i)=3u^2/4,\qquad E_6(i)=0.
\]
The rational modular-form ring is $\Q[E_4,E_6]$
\cite[Proposition 4]{Z}. Thus $f_0=\alpha u^{2p}$ for some
$\alpha\in\Q$. For a weight-$k$ modular form $F$, its Serre derivative
\[
 \vartheta_kF=DF-\frac{k}{12}E_2F
\]
is modular of weight $k+2$, with rational Fourier coefficients
when $F$ has them \cite[Section 5.1]{Z}. Put
$F=\vartheta_w f$ and $H=\vartheta_{w+2}F$.
Since a weight congruent to two modulo four vanishes at $i$,
$F(i)=0$. Ramanujan's identity
$DE_2=(E_2^2-E_4)/12$ therefore gives
\begin{align*}
 f_2
 &=\frac{w}{12}(DE_2)(i)f_0
       +\frac{w}{12}E_2(i)f_1+H(i)\\
 &=\frac{w(w+1)}{16\pi^2}f_0
                 +H(i)-\frac{w}{144}E_4(i)f_0.
\end{align*}
The last two terms are the value at $i$ of a rational modular form
of weight $w+4$, so they equal $\beta u^{2p+2}$ for some
$\beta\in\Q$. Hence
\begin{equation}\label{split}
 f_0=\alpha u^{2p},\qquad f_1=\frac{w}{4\pi}f_0,
 \qquad f_2=\frac{w(w+1)}{16\pi^2}f_0+\beta u^{2p+2}.
\end{equation}

Put $L=w(w+1)f_0/(8\pi)$ and $U=-2\pi\beta u^{2p+2}$.
Substitution of \eqref{split} in \eqref{Iminus}--\eqref{Iplus} gives
\[
 I_p^-=4ca\pi^{w+1}(L+U),\qquad
 I_p^+=-4ca\pi^{w+1}\{(1-t)L+(1+t)U\}.
\]
The first part is a rational multiple of $G^{8p}/\pi^{2p}$,
and the second is a rational multiple of $G^{8p+8}/\pi^{2p+4}$.
These two numbers are linearly independent over $\Q$, by the
established algebraic independence of $\pi$ and $\Gamma(1/4)$
\cite[p.~85]{Z}. Comparison with \eqref{minus}--\eqref{plus}
therefore proves $g_p=-(1-t)i_p$ and $h_p=-(1+t)j_p$, as required.
All steps apply to every $p\ge1$.\qed

For an explicit coefficient algorithm, the same calculation gives
\[
 i_p=\frac{(-1)^p c w(w+1)\alpha}{2^{2p+2}},\qquad
 j_p=\frac{(-1)^{p+1}c\beta}{2^{2p+2}}.
\]
The quantities $\alpha,\beta$ follow by rational polynomial
arithmetic in $E_4,E_6$ and the three Ramanujan differential equations.
For example, $p=1$ gives $c=1/240$, $\alpha=3/4$, $\beta=5/64$, and
\[
 (i_1,j_1,g_1,h_1)=
 \left(-\frac1{256},\frac1{3\cdot2^{14}},
              \frac3{512},-\frac1{3\cdot2^{15}}\right),
\]
in agreement with the source examples. This checks the smallest
allowed index without relying on an asymptotic or generic-index argument.

\section*{Task 2. Generalize the modular correspondence}
Replace the degree-two selector identity by the degree-$N$ Eisenstein
identity. Identify integral families whose residue sums realize the
resulting transformed derivatives, and derive their coefficient relations.

\section*{Task 3. Compare higher denominator powers}
Find the analogue of the two coefficient relations when the denominator
power is four or higher. Determine which finite combinations of modular
derivatives control each period power, including contributions from
the additional poles and the convergence condition at zero.

\section*{Task 4. Determine coefficient signs and growth}
Determine the signs and possible vanishing of $\beta$ in \eqref{split}
for all $p$. Obtain asymptotic formulas for $i_p,j_p$ with explicit error
bounds, and use the proved relations to transfer them to $g_p,h_p$.

\section*{Task 5. Derive recurrences for exact coefficients}
Construct recurrences for $\alpha$ and $\beta$ as the weight increases
by four. Seek a method computing the four integral coefficients without
reconstructing the full Eisenstein polynomial at each new weight.

\begin{thebibliography}{2}
\bibitem{XZ} C.~Xu and J.~Zhao,
\emph{Berndt-Type Integrals and Series Associated with Ramanujan and
Jacobi Elliptic Functions},
\href{https://arxiv.org/abs/2301.08211v2}{arXiv:2301.08211v2} (2024),
Theorems 6.3 and 8.12, Conjecture 8.14.
\bibitem{Z} D.~Zagier, \emph{Elliptic Modular Forms and Their Applications},
in \emph{The 1-2-3 of Modular Forms}, Springer (2008), pp.~1--103;
\href{https://people.mpim-bonn.mpg.de/zagier/files/doi/10.1007/978-3-540-74119-0_1/fulltext.pdf}{author-hosted chapter}.
\end{thebibliography}
\noindent\textit{Verification and status.}
Run \texttt{verify\_relations.py} for exact symbolic checks and separate
numerical diagnostics. The proof for all $p$ is the analytic and
modular argument above, not a finite test. Full current-source and
dated bounded claim checks accompany this note. The authoring agent
conducted the proof audit; independent peer review has not occurred.
\end{document}
